\let\oepi\epigraph
\renewcommand{\epigraph}[2]{\epigraphlist{#1}{#2}\oepi{#1}{#2}}

\makeatletter
\newcommand{\openingquote}[2]{\gdef\@openquote{#1}\gdef\@openauthor{#2}}
\newcommand{\openingquotealign}{center}

\newcommand{\placeopeningquote}{%
\newpage
\begingroup
    \vspace*{\fill}
    \begin{\openingquotealign}
        \begin{tikzpicture}[every node/.style={inner sep=0pt}]
            \node[text width=\textwidth/1.3,align=center](Text){%
                    \fontfamily{Crimson-TLF}\selectfont
                    \vspace{0.2in}
                    \LARGE
                    {\\\@openquote \\}
                    \vspace{0.1in}
                    \pgfornament[width=0.3\textwidth]{89} \\
                    \vspace{0.2in}
                    {\large\itshape\@openauthor}
                    \vspace{0.3in}
                    } ;
            \node[shift={(-1cm,1cm)},anchor=north west](CNW)
            at (Text.north west) {\pgfornament[width=1.75cm]{61}};
            \node[shift={(1cm,1cm)},anchor=north east](CNE)
            at (Text.north east) {\pgfornament[width=1.75cm,symmetry=v]{61}};
            \node[shift={(-1cm,-1cm)},anchor=south west](CSW)
            at (Text.south west) {\pgfornament[width=1.75cm,symmetry=h]{61}};
            \node[shift={(1cm,-1cm)},anchor=south east](CSE)
            at (Text.south east) {\pgfornament[width=1.75cm,symmetry=c]{61}};
            \pgfornamenthline{CNW}{CNE}{north}{84}
            \pgfornamenthline{CSW}{CSE}{south}{84}
            \pgfornamentvline{CNW}{CSW}{west}{84}
            \pgfornamentvline{CNE}{CSE}{east}{84}
        \end{tikzpicture}
    \end{\openingquotealign}
    \vspace*{\fill}
\endgroup}

\newcommand{\closingquote}[2]{\gdef\@closequote{#1}\gdef\@closeauthor{#2}}
\newcommand{\closingquotealign}{center}
\newcommand{\placeclosingquote}{%
\newpage
\fancyhead[LE,LO]{}
\fancyhead[RE,RO]{}
\renewcommand{\headrulewidth}{0pt}
\begingroup
    \vspace*{\fill}
    \begin{\closingquotealign}
        \begin{tikzpicture}[every node/.style={inner sep=0pt}]
            \node[text width=\textwidth/1.3,align=center](Text){%
                    \fontfamily{Crimson-TLF}\selectfont
                    \vspace{0.1in}
                    \LARGE
                    {\\\@closequote \\}
                    \vspace{0.1in}
                    \pgfornament[width=0.3\textwidth]{89} \\
                    \vspace{0.1in}
                    {\large\itshape\@closeauthor}
                    \vspace{0.3in}
                    } ;
            \node[shift={(-1cm,1cm)},anchor=north west](CNW)
            at (Text.north west) {\pgfornament[width=1.75cm]{61}};
            \node[shift={(1cm,1cm)},anchor=north east](CNE)
            at (Text.north east) {\pgfornament[width=1.75cm,symmetry=v]{61}};
            \node[shift={(-1cm,-1cm)},anchor=south west](CSW)
            at (Text.south west) {\pgfornament[width=1.75cm,symmetry=h]{61}};
            \node[shift={(1cm,-1cm)},anchor=south east](CSE)
            at (Text.south east) {\pgfornament[width=1.75cm,symmetry=c]{61}};
            \pgfornamenthline{CNW}{CNE}{north}{84}
            \pgfornamenthline{CSW}{CSE}{south}{84}
            \pgfornamentvline{CNW}{CSW}{west}{87}
            \pgfornamentvline{CNE}{CSE}{east}{87}
        \end{tikzpicture}
  \end{\closingquotealign}
  \vspace*{\fill}
\endgroup}


\newcommand{\dedication}[1]{\gdef\@dedicate{#1}}
\newcommand{\dedicationalign}{center}
\newcommand{\placededication}{%
\newpage
\begingroup
    \vspace*{\fill}
    \begin{\dedicationalign}
        \begin{tikzpicture}[every node/.style={inner sep=0pt}]
            \node[text width=\textwidth/1.3,align=center](Text){%
                    \fontfamily{Crimson-TLF}\selectfont
                    \vspace{0.2in}
                    \LARGE
                    {\\\@dedicate \\}
                    \phantom{ }
                    \vspace{0.3in}
                    } ;
            \node[shift={(-1cm,1cm)},anchor=north west](CNW)
            at (Text.north west) {\pgfornament[width=1.75cm]{61}};
            \node[shift={(1cm,1cm)},anchor=north east](CNE)
            at (Text.north east) {\pgfornament[width=1.75cm,symmetry=v]{61}};
            \node[shift={(-1cm,-1cm)},anchor=south west](CSW)
            at (Text.south west) {\pgfornament[width=1.75cm,symmetry=h]{61}};
            \node[shift={(1cm,-1cm)},anchor=south east](CSE)
            at (Text.south east) {\pgfornament[width=1.75cm,symmetry=c]{61}};
            \pgfornamenthline{CNW}{CNE}{north}{84}
            \pgfornamenthline{CSW}{CSE}{south}{84}
            \pgfornamentvline{CNW}{CSW}{west}{84}
            \pgfornamentvline{CNE}{CSE}{east}{84}
        \end{tikzpicture}
    \end{\dedicationalign}
    \vspace*{\fill}
\endgroup}

\renewcommand{\epigraphsize}{\small}
\setlength{\epigraphwidth}{0.65\textwidth}
\renewcommand{\textflush}{flushright}
\renewcommand{\sourceflush}{flushright}
\renewcommand{\epigraphflush}{flushright}
% A useful addition
\newcommand{\epitextfont}{}
\newcommand{\episourcefont}{\scshape}

\newsavebox{\epi@textbox}
\newsavebox{\epi@sourcebox}
\newlength\epi@finalwidth
\renewcommand{\epigraph}[2]{%
  \vspace{\beforeepigraphskip}
  {\epigraphsize\begin{\epigraphflush}
  \epi@finalwidth=\z@
  \sbox\epi@textbox{%
     \varwidth{\epigraphwidth}
     \begin{\textflush}\epitextfont#1\end{\textflush}
     \endvarwidth
  }%
  \epi@finalwidth=\wd\epi@textbox
  \sbox\epi@sourcebox{%
     \varwidth{\epigraphwidth}
     \begin{\sourceflush}\episourcefont#2\end{\sourceflush}%
     \endvarwidth
  }%
  \ifdim\wd\epi@sourcebox>\epi@finalwidth
     \epi@finalwidth=\wd\epi@sourcebox
  \fi
  \leavevmode\vbox{
     \hb@xt@\epi@finalwidth{\hfil\box\epi@textbox}
     \vskip1.75ex
     \hrule height \epigraphrule
     \vskip.75ex
     \hb@xt@\epi@finalwidth{\hfil\box\epi@sourcebox}
  }%
  \end{\epigraphflush}
  \vspace{\afterepigraphskip}}}


\newcommand{\inplaceEpigraph}[4][\fill]{%
  \vspace{#1}
  {\epigraphsize\begin{center}
     \begin{center}\epitextfont#2\end{center}
     \vspace{-5px}
     \rule{#4\textwidth}{.4pt}
     \vspace{-10px}
     \begin{center}\episourcefont#3\end{center}%
      \end{center}
  }%
  \vspace{#1}
  }


\newcommand{\mediaepigraph}[4][by]{\epigraph{#2}{\textit{#3}, #1\ #4}}
\newcommand{\showepigraph}[3]{\epigraph{#1}{\textit{#2}, ``#3''}}

\newcommand{\translationEpigraph}[3]{%
  \vspace{\beforeepigraphskip}
  \vbox{%
    \setlength{\epigraphwidth}{0.48\textwidth}
    \xpatchcmd{\@epitext}{{minipage}}{{minipage}[t]}{}{}%
    \xpatchcmd{\@epitext}{1\\}{1}{}{}%
    \epigraphsize
    \begin{\epigraphflush}
    \begin{minipage}[b]{0.48\textwidth}
      \@epitext{\epitextfont#1}
      \@episource{\episourcefont#3\strut}%
    \end{minipage}\hfill\vline\hfill
    \begin{minipage}[b]{0.48\textwidth}
      \@epitext{\begin{flushleft}\epitextfont#2\end{flushleft}}
      \@episource{\begin{flushleft}\episourcefont#3\strut\end{flushleft}}%
    \end{minipage}%
    \end{\epigraphflush}%
  }%
  \nointerlineskip
  \vspace*{\afterepigraphskip}%
  \ifepigraphnoindent\@afterheading\fi
}

\newcommand{\twinEpigraphs}[4]{%
  \vspace*{\fill}
  \noindent
  \begin{tcbraster}[raster columns=2,raster equal height=rows,valign=bottom,enhanced,opacityfill=0,frame hidden]
    \begin{tcolorbox}
    {\epigraphsize\begin{center}
     \begin{center}\epitextfont#1\end{center}%
     % \vspace{-5px}
     \rule{\textwidth}{.4pt}
     % \vspace{-20px}
     \begin{center}\episourcefont#2\end{center}%
    \end{center}
    }
    \end{tcolorbox}
    \begin{tcolorbox}
    {\epigraphsize\begin{center}
     \begin{center}\epitextfont#3\end{center}%
     % \vspace{-5px}
     \rule{\textwidth}{.4pt}
     % \vspace{-20px}
     \begin{center}\episourcefont#4\end{center}%
    \end{center}
    }
    \end{tcolorbox}
  \end{tcbraster}
  \vspace*{\fill}
}



\newcommand{\fancifulepigraph}[4]{%
  \begin{fanciful}
  \vspace{\beforeepigraphskip}
  \vbox{%
    \setlength{\epigraphwidth}{0.48\textwidth}
    \xpatchcmd{\@epitext}{{minipage}}{{minipage}[t]}{}{}%
    \xpatchcmd{\@epitext}{1\\}{1}{}{}%
    \epigraphsize
    \begin{\epigraphflush}
    \begin{minipage}[b]{0.48\textwidth}
      \@epitext{\epitextfont#1}
      \@episource{\citet{#3}\strut}%
    \end{minipage}\hfill\vline\hfill
    \begin{minipage}[b]{0.48\textwidth}
      \@epitext{\begin{flushleft}\epitextfont#2\end{flushleft}}
      \@episource{\begin{flushleft}\citet{#4}\strut\end{flushleft}}%
    \end{minipage}%
    \end{\epigraphflush}%
  }%
  % \nointerlineskip
  % \vspace*{\afterepigraphskip}%
  % \ifepigraphnoindent\@afterheading\fi
  \end{fanciful}
}

\usepackage{fmtcount}
\usepackage{listofitems}

\setsepchar{?}
\readlist\alice@list{Six impossible things. Count them, Alice.?There's a potion that can make you shrink.?And a cake that can make you grow.?Animals can talk.?Cats can disappear.?There's a place called Wonderland.?I can slay the Jabberwocky.}

\newcounter{impossible}
\setcounter{impossible}{-1}
\newcommand{\alice@phrase}{}
\newcommand{\alice@attr}{}
\newcommand{\alice}{%
    \stepcounter{impossible}
    \ifnum\theimpossible=\numexpr\alice@listlen-1
        \def\alice@ordinal{Last}
    \else
        \def\alice@ordinal{\Ordinalstring{impossible}}
    \fi

    \ifnum\theimpossible=0
        \renewcommand{\alice@attr}{Alice in Wonderland}
    \else
        \renewcommand{\alice@attr}{\alice@ordinal~of the Six Impossible Things}
    \fi
    \renewcommand{\alice@phrase}{\alice@list[\theimpossible+1]}
    \epigraph{\alice@phrase}{\alice@attr}
}
\makeatother

\newcommand{\chapterEpigraph}[2]{\tikz[remember picture,overlay]
    {
    \node[fill=main-accent!50!white,font=\fontsize{10}{12}\selectfont\color{white},anchor=west, minimum height=\ChapWd, inner sep=0.5cm]
      at ([xshift=60pt,yshift=-60pt-0.5\ChapWd]current page.north west)
      (epigraphloc) {\begin{varwidth}{0.5\textwidth}#1\begin{flushright}\scshape-- #2\end{flushright}\end{varwidth}};
    }}